\section{Conclusion}

This paper presents a controlled Isaac Sim boundary study of how
planner-to-executor contract design and runtime validation shape
embodied-agent execution. In the fixed task matrix studied here, under-specified direct-action handoffs fail at the
planner/executor boundary, typed tool calls remove that failure mode in the
fixed-R0 slice, and validate-and-retry recovers the recoverable invalid-first-action
cases.

The failure analysis shows that the dominant errors are contract violations,
not generic task collapse. Explicit contracts prevent those failures at
the source, runtime validation repairs some that remain, and a brief
self-check can reduce planner/tool call counts within an already strong contract.
Because the planner backend is deterministic, these are descriptive
task-instance comparisons rather than stochastic rollout estimates. Within the
tested Isaac Sim setting, the practical implication is to make the callable
interface explicit, validate it before dispatch, and treat planner-side
self-checks as an efficiency-oriented refinement rather than a separate
success mechanism.
